\documentclass[10pt,a4paper]{amsart}
\usepackage[margin=19mm]{geometry}
\usepackage{amsmath,amssymb,mathtools}
\usepackage[scr=rsfs,cal=euler]{mathalfa}
\usepackage{xfrac}
\usepackage[unicode,hidelinks,bookmarks=false]{hyperref}
\newcommand{\ie}{i.e.\ }
\newcommand{\cf}{cf.\ }
\newcommand{\resp}{respectively\ }
\newcommand{\Subsection}{\subsection}
\providecommand{\nobreakdash}{\nobreak\hbox{-}}
\setlength{\emergencystretch}{2em}
\multlinegap0pt
\DeclareMathOperator{\GL}{GL}
\usepackage[shortlabels]{enumitem}
\usepackage{mathtools}
\usepackage{amssymb}
\usepackage{xcolor}
\newtheorem{lemma}{Lemma}
\newcommand{\VV}{\mathcal{V}}
\newcommand{\Z}{\mathbb{Z}}
\title{Strong Hybrid Subconvexity for Twisted Selfdual $\GL_3$ $L$-Functions}
\author{Soumendra Ganguly, Peter Humphries, Yongxiao Lin,, Ramon Nunes}
\date{Source: arXiv:2408.00596v2 (CC BY 4.0)}
\begin{document}
\maketitle

\noindent\emph{Source and licence.} Strong Hybrid Subconvexity for Twisted Selfdual $\GL_3$ $L$-Functions. Version arXiv:2408.00596v2 (CC BY 4.0), distributed by arXiv under the Creative Commons Attribution 4.0 International licence, http://creativecommons.org/licenses/by/4.0/, as recorded in the arXiv licence metadata for this exact version and shown on the abstract page https://arxiv.org/abs/2408.00596v2. Full licence notice: \texttt{licenses/arxiv-2408.00596-CC-BY-4.0.txt}.\par\smallskip

\noindent\textit{Editorial scope.} Counted unit: the article's lemma lem:VVchiprincprinc (source0.tex lines 685--698). The article's complete original proof of it is delivered with it (source0.tex lines 700--735, one \texttt{proof} environment of 36 lines). No mathematics is written, repaired or re-derived: the statement and the proof are the article's own lines, in the article's own notation, and no retained line is substituted or re-flowed.  Retained uncounted as the source-local context the proof needs: lines 452--455.  Nothing was omitted from any retained span, so the package carries no omission record.  No retained line contains a citation, so the wrapper carries no bibliography and no bibliography entry.  No retained line contains a figure environment, an image-inclusion command or drawing code, so no image is delivered and the package declares no asset dependency.  Editorial formatting only, in addition: an AMS \texttt{amsart} preamble that restates the article's own declarations and macros used by the retained lines, namely the environments \texttt{lemma} on separate counters, together with the standard packages those lines need.  The retained environments print in this document as Lemma 1 in order of appearance, whereas the article prints the same items with the numbering of its own full document.  The complete native source of the article as distributed in the arXiv e-print for this exact version is bundled unchanged as \texttt{sources/163213/source0.tex}.
\begin{equation}
\label{eqn:Rqndefeq}
R_q(n) \coloneqq \sum_{d \in (\Z/q\Z)^{\times}} e\left(\frac{dn}{q}\right) = \sum_{d \mid (q,n)} d \mu\left(\frac{q}{d}\right)
\end{equation}

\begin{lemma}
\label{lem:VVchiprincprinc}
Let $\chi_{0(p^{\beta})},\psi_{0(p^{\beta})}$ both be the principal character modulo $p^{\beta}$. Suppose that $m_3 \mid m_1$, $m_1 \mid m_2$, and $m_1,m_2,m_3,r \mid p^{\infty}$. Then
\[\VV_{\chi_{0(p^{\beta})}}(\psi_{0(p^{\beta})};m_1,m_2,m_3,r) = \begin{dcases*}
p^{-1} (p - 1)^3 & if $r = 1$, $p \mid m_1,m_2,m_3$, and $\beta = 1$,	\\
-p^{-1} (p - 1)^2 & if $m_3, r = 1$, $p \mid m_1,m_2$, and $\beta = 1$,	\\
p^{-1} (p - 1) & if $m_1, m_3, r = 1$, $p \mid m_2$, and $\beta = 1$,	\\
p^{2\beta - 1} (p - 1) & if $m_1, m_2, m_3 = 1$ and $p \mid r$,	\\
p^{-1} (p^3 - p^2 - p - 1) & if $m_1, m_2, m_3, r = 1$ and $\beta = 1$,	\\
p^{2\beta - 1}(p - 1) & if $m_1, m_2, m_3, r = 1$ and $\beta \geq 2$,	\\
0 & otherwise.
\end{dcases*}\]
In particular, $\VV_{\chi_{0(p^{\beta})}}(\psi_{0(p^{\beta})};m_1,m_2,m_3,r) \ll p^{2\beta}$.
\end{lemma}

\begin{proof}
The character sum of interest is
\begin{multline}
\label{eqn:VVchiprincprincdefeq}
\VV_{\chi_{0(p^{\beta})}}(\psi_{0(p^{\beta})};m_1,m_2,m_3,r)	\\
= \frac{1}{p^{\beta}} \sum_{t,u \in \Z/p^{\beta}\Z} R_{p^{\beta}}(t + m_2 u) \chi_{0(p^{\beta})}(rt + m_1 m_2) R_{p^{\beta}}(u) \chi_{0(p^{\beta})}(ru - m_1) R_{p^{\beta}}(m_3 t),
\end{multline}
where $R_q(n)$ denotes the Ramanujan sum, as in \eqref{eqn:Rqndefeq}. To determine its exact value, we must treat this on a case-by-case basis. Below, we freely use the fact that
\begin{equation}
\label{eqn:Rpbetapalphadefeq}
R_{p^{\beta}}(p^{\alpha}) = \begin{dcases*}
0 & if $0 \leq \alpha \leq \beta - 2$,	\\
-p^{\beta - 1} & if $\alpha = \beta - 1$,	\\
p^{\beta - 1}(p - 1) & if $\alpha \geq \beta$.
\end{dcases*}
\end{equation}
\begin{itemize}[leftmargin=*]
\item If $m_1,m_2,m_3 \equiv 0 \pmod{p}$, then the summand in \eqref{eqn:VVchiprincprincdefeq} vanishes unless $r = 1$ and $(tu,p) = 1$, in which case it is
\[\frac{1}{p^{\beta}} \sum_{t,u \in (\Z/p^{\beta}\Z)^{\times}} R_{p^{\beta}}(t + m_2 u) R_{p^{\beta}}(u) R_{p^{\beta}}(m_3 t).\]
The second term in the sum vanishes unless $u \equiv 0 \pmod{p^{\beta - 1}}$, which can only occur if $\beta = 1$. In this case, the first term and second terms are both $-1$, while the third is $p - 1$, and so we obtain $p^{-1} (p - 1)^3$.
\item If $m_1,m_2 \equiv 0 \pmod{p}$ and $m_3 = 1$, then we follow the same argument above except the third term in the sum is $-1$, and so we obtain $-p^{-1} (p - 1)^2$.
\item If $m_1, m_3 = 1$ and $m_2 \equiv 0 \pmod{p}$, then the summand vanishes unless $r = 1$ and $(t,p) = 1$. We make the change of variables $t \mapsto \overline{r} t$ and $u \mapsto \overline{r} (u + 1)$, yielding
\[\frac{1}{p^{\beta}} \sum_{t,u \in (\Z/p^{\beta}\Z)^{\times}} R_{p^{\beta}}(t + m_2 u + m_2) R_{p^{\beta}}(u + 1) R_{p^{\beta}}(t).\]
The third term in the sum vanishes unless $\beta = 1$, in which case the first and third terms are $-1$, while the second is $-1$ unless $u = p - 1$, in which case it is $p - 1$, and so we obtain $p^{-1} (p - 1)$.
\item If $m_1, m_2, m_3 = 1$ and $r \equiv 0 \pmod{p}$, this is
\[\frac{1}{p^{\beta}} \sum_{t,u \in \Z/p^{\beta}\Z} R_{p^{\beta}}(t - u) R_{p^{\beta}}(u) R_{p^{\beta}}(t).\]
The summand vanishes unless $t,u \equiv 0 \pmod{p^{\beta - 1}}$. Making the change of variables $t \mapsto p^{\beta - 1}t$ and $u \mapsto p^{\beta - 1}u$, this becomes
\[p^{2\beta - 3} \sum_{t,u \in \Z/p\Z} R_p(t - u) R_p(u) R_p(t),\]
which is $p^{2\beta - 1} (p - 1)$.
\item If $m_1, m_2, m_3, r = 1$, then we instead make the change of variables $t \mapsto t - 1$ and $u \mapsto 1 - u$, yielding
\[\frac{1}{p^{\beta}} \sum_{t,u \in (\Z/p^{\beta}\Z)^{\times}} R_{p^{\beta}}(t - u) R_{p^{\beta}}(u - 1) R_{p^{\beta}}(t - 1).\]
For $\beta = 1$, this is $p^{-1} (p^3 - p^2 - p - 1)$. For $\beta \geq 2$, the summand vanishes unless $t,u \equiv 1 \pmod{p^{\beta - 1}}$. Making the change of variables $t \mapsto 1 + p^{\beta - 1}t$ and $u \mapsto 1 + p^{\beta - 1}u$, this becomes
\[p^{2\beta - 3} \sum_{t,u \in \Z/p\Z} R_p(t - u) R_p(u) R_p(t),\]
which is again $p^{2\beta - 1} (p - 1)$.\qedhere
\end{itemize}
\end{proof}

\end{document}
